% !TeX root = main.tex
% !TeX spellcheck = pt_BR

\chapter{Álgebra Linear}

\section{Sistemas de equações lineares}

\begin{theorem}[Sistemas de equações lineares]
	 Todo sistema de equações lineares tem ou nenhuma solução, ou exatamente uma solução ou infinitas soluções.
	 \begin{note}
	 	Um sistema que não possui solução é chamado de \underline{inconsistente}; se existir pelo menos uma solução, dizemos que o sistema é \underline{consistente}.
	 \end{note}
\end{theorem}

\begin{theorem}[Sistemas homogêneos]
	\leavevmode
	\begin{outline}[enumerate]
		\1 Um sistema homogêneo de equações lineares com mais incógnitas que equações (subdeterminado) tem infinitas soluções.
		\1 Se o sistema homogêneo tem $n$ incógnitas e a forma escolanada reduzida por linhas (echelon) da matriz aumentada tem $r$ linhas não nulas, então o sistema possui $n-r$ variáveis livres (parâmetros).
	\end{outline}
		\begin{note}
		\leavevmode
		\begin{enumerate}
			\item Um sistema não-homogêneo com mais equações que incógnitas (sobredeterminado) não precisa ser consistente.
			\item Todo sistema homogêneo é consistente, pois admite ao menos a solução trivial.
		\end{enumerate}
	\end{note}
\end{theorem}

\begin{theorem}[Solução geral e particular] Seja $\v{x}_0$ qualquer solução de um sistema consistente $\m{A}\v{x}=\v{b}$. Se $S=\{\v{v}_1, \v{v}_2, \ldots, \v{v}_k\}$ é uma base para o espaço-nulo de $\m{A}$, então toda solução de $\m{A}\v{x}=\v{b}$ pode ser expressa como
	\begin{equation}
		\v{x} = \v{x}_0 + c_1\v{v}_1 + c_2\v{v}_2 + \cdots + c_k\v{v}_k.
	\end{equation}
	De forma contrária, para todas as escolhas de $c_1. c_2, \ldots, c_k$, $\v{x}$ é uma solução de $\m{A}\v{x}=\v{b}$.
	\begin{note}
		$\v{x}_0$ é chamado de solução particular de $\m{A}\v{x}=\v{b}$ e a parte restante de solução geral de $\m{A}\v{x}=\v{0}$. Assim, a solução geral de um sistema consistente de equações lineares pode ser expressa como a soma de uma solução particular do sistema e a solução geral do sistema homogêneo correspondente.
	\end{note}
\end{theorem}

\begin{theorem}[Sistema $m \times n$]. Seja $\m{A}_{m\times n}$ com $\pos(\m{A})=r$.
	\begin{outline}
		\1 \textbf{Caso sobredeterminado}: se $m>n$, então $\exists\, \v{b} \in \Re^m$ tal que o sistema $\m{A}\v{x}=\v{b}$ é inconsistente (sempre, pois $r \leq n < m \implies \Col(\m{A}) \subsetneq \Re^m$).
		\2 Para os $\v{b} \in \Col(\m{A})$ (onde o sistema \emph{é} consistente):
		\3 Se $r=n$ (posto-coluna completo): solução única ($n-r=0$ parâmetros);
		\3 Se $r<n$: solução geral com $n-r>0$ parâmetros (infinitas soluções).
		\1 \textbf{Caso subdeterminado}: se $m<n$, então $\forall\, \v{b} \in \Re^m$, o sistema $\m{A}\v{x}=\v{b}$ é inconsistente ou possui infinitas soluções (nunca única, pois $r \leq m < n$).
		\2 Se $r=m$ (posto-linha completo): $\Col(\m{A})=\Re^m$, logo o sistema é \emph{sempre} consistente, com $n-r>0$ parâmetros;
		\2 Se $r<m$: o sistema é consistente sse $\v{b} \in \Col(\m{A})$; se consistente, possui $n-r$ parâmetros.
	\end{outline}
\end{theorem}

\section{Definições e propriedades de matrizes}

\begin{definition} Definições de matrizes
	\leavevmode
	\begin{outline}[enumerate]
		\1 Matriz singular: Se $\m{A}$ não admite inversa, $\m{A}$ é dita singular.
		\1 Matriz simétrica: $\m{A} = \m{A}\T$
		\1 Matriz ortogonal: \( \m{A}\inv = \m{A}\T \)
		\1 Matriz triangular superior (inferior): elementos abaixo (acima) da diagonal principal são nulos.
		\1 Matriz quadrada idempotente: $\m{A}$ é idempotente se $A=AA$
		\1 Traço: $\tr{\m{A}} = a_{11} + a_{22} + \cdots + a_{nn}$
		\1 \( \m{A}^0 = \m{I} \)
		\1 \( \m{A}^n = \underbrace{\m{A}\m{A}\cdots\m{A}}_{\text{n fatores}} \)
		\1 \( \m{A}^{-n} = (\m{A}\inv)^n \) e
		\1 Polinômio matricial: Seja $\m{A}$ uma matriz quadrada e $p(x) = a_0 + a_1 x + \cdots + a_n x^n$. Definimos $p(\m{A}) = a_0\m{I} + a_1\m{A} + \cdots + a_n\m{A}^n$.
		\1 $\v{u}$ e $\v{v}$ são ortogonais se $\v{u}\cdot\v{v} = \v{u}\T\v{v} = 0$
	\end{outline}
\end{definition}

\begin{theorem}[Lei dos expoentes]
	\leavevmode
	\begin{outline}[enumerate]
		\1 \( \m{A}^r\m{A}^s = \m{A}^{r+s} \)
		\1 \( (\m{A}^r)^s = A^{rs} \)
	\end{outline}
\end{theorem}

\begin{theorem}[Propriedades da transposta]
	\leavevmode
	\begin{outline}[enumerate]
		\1 \( (\m{A}\T)\T = \m{A} \)
		\1 \( (\m{A}+\m{B})\T = \m{A}\T + \m{B}\T \) 
		\1 \( (\m{A}-\m{B})\T = \m{A}\T - \m{B}\T \)
		\1 \( (k\m{A})\T = k\m{A}\T \)
		\1 \( (\m{A}\m{B})\T = \m{B}\T \m{A}\T \) 
	\end{outline}
\end{theorem}


\begin{theorem}[Propriedades de matrizes inversas]
	\leavevmode
	\begin{outline}[enumerate]
		\1 Se $\m{A}$ admite inversa, então a inversa é única
		\1 \( (\m{A}\inv)\inv = \m{A} \)
		\1 \( \m{A} \text{ é invertível} \implies \m{A}^n \text{ é invertível e } (\m{A}^n)\inv = (\m{A}\inv)^n \)
		\1 \( \m{A} \text{ é invertível} \iff \m{A}\T\m{A} \text{ e } \m{A}\m{A}\T \text{ são invertíveis} \)
		\1 \( (\m{A}\m{B})\inv = \m{B}\inv\m{A}\inv \)
		\1 \( (k\m{A})\inv = \frac{1}{k}\m{A}\inv \)
		\1 \( (\m{A}\T)\inv = (\m{A}\inv)\T \)
		\1 Sejam $\m{A}$ e $\m{B}$ são matrizes quadradas de mesma dimensão. \\
		$\m{A}\m{B}$ é invertível $\iff$ $\m{A}$ e $\m{B}$ também são.
	\end{outline}
\end{theorem}

\begin{theorem}[Inversa de matriz $2 \times 2$]
	Seja $\m{A}_{2\times 2}$ definida por
	\begin{equation}
		\m{A} = \begin{bmatrix*}[r]
			a & b \\
			c & d
		\end{bmatrix*}
	\end{equation}
	Se $ad-bc \neq 0$, a inversa de $\m{A}$ é dada por
	\begin{equation}
		\m{A}\inv = \frac{1}{ad-bc} 
		\begin{bmatrix*}[r]
			d & -b \\
			-c & a
		\end{bmatrix*}
	\end{equation}
\end{theorem}

\begin{theorem}[Propriedades de matrizes triangulares]
	\leavevmode
	\begin{outline}[enumerate]
		\1 A transposta de uma matriz triangular inferior (superior) é triangular inferior (superior).
		\1 O produto de matrizes triangulares inferiores (superiores) é triangular inferior (superior).
		\1 Uma matriz triangular é invertível se, e somente se, as entradas na diagonal principal são todas não nulas.
		\1 A inversa de uma matriz triangular inferior (superior) é triangular inferior (superior).
		\1 $\m{A}$ é triangular $\implies$ $\det(\m{A})=a_{11} a_{22} \cdots a_{nn}$
		\1 $\m{A}$ é triangular $\implies$ os autovalores de $\m{A}$ são as entradas da diagonal principal
		
	\end{outline}
\end{theorem}

\begin{theorem}[Propriedades de matrizes simétricas]
	\leavevmode
	\begin{outline}[enumerate]
		\1 $\m{A}\T \m{A}$ e $\m{A} \m{A}\T$ são simétricas
		\1 \( \m{A}, \m{B} \text{ são simétricas } \implies \m{A}+\m{B} \text{ e } \m{A}-\m{B}
		 \text{ são simétricas} \)
		 \1 \( \m{A} \text{ é simétrica} \implies \m{A}\inv \text{ é simétrica } \)
		 \1 \( \m{A} \text{ é invertível} \implies \m{A}\m{A}\T \text{ e }  \m{A}\T\m{A} \text{ são invertíveis.}  \)
	\end{outline}
\end{theorem}

\begin{theorem}[Propriedades de matrizes idempotentes]
	\leavevmode
	\begin{outline}[enumerate]
		\1 $\m{A}$ é idempotente $\implies$ $\pos(\m{A}) = \tr(\m{A})$
		\1 $\m{A}-\m{B}$ é idempotente $\implies$ $\pos(\m{A}-\m{B}) = \pos(\m{A}) - \pos(\m{B})$
	\end{outline}
\end{theorem}

\begin{theorem}[Propriedades da matriz Design e Matriz $\hat{H}$]
	Sejam:
	\begin{itemize}
		\item $\m{X}_{n \times (p+1)}$ com $n > p+1$ e posto $p+1$ (posto-coluna) completo;
		\item $\hat{\m{H}}_{n \times n} = \m{X}(\m{X}\T\m{X})\inv\m{X}\T$.
	\end{itemize} Então:
	\begin{outline}[enumerate]
		\1 $\m{H}$ é idempotente.
		\1 $\m{I}-\m{H}$ é idempotente.
		\1 $\m{H}$ e $\m{I}-\m{H}$ são ortogonais, isto é, $(\m{I}-\m{H})\m{H} = \m{0}$.
		\1 \( \pos(\m{H}) = \pos(\m{X}) = p+1\)
		\1 \(\pos(\m{I}-\m{H}) =n-(p+1) \)
	\end{outline}
\end{theorem}

\begin{theorem}[Propriedades do traço]
		\leavevmode
	\begin{outline}[enumerate]
		\1 $\tr(\m{A}\m{B}) = \tr(\m{B}\m{A})$
		\1 \( \tr(\m{A} + \m{B}) = \tr(\m{A}) + \tr(\m{B}) \)
	\end{outline}
\end{theorem}

\begin{theorem}[Propriedades do determinante]
	\leavevmode
	\begin{outline}[enumerate]
		\1 \( \det(\m{A}) = \det(\m{A}\T) \)
		\1 \( \det{(\m{A}\m{B})} = \det{\m{A}}\det{\m{B}} \)
		\1 \( \det(k\m{A}) = k^n \det(\m{A})\)
		\1 \( \m{A} \text{ é invertível} \iff  \det{\m{A}} \neq 0 \)
		\1 \( \m{A} \text{ é invertível} \implies  \det(\m{A}\inv) = \frac{1}{\det\m{A}}\)
		\1 Se $\m{A}$ tem uma linha (coluna) de zeros, então $\det(\m{A})=0$
		\1 Se $\m{A}$ tem duas linhas (colunas) proporcionais, então $\det(\m{A})=0$
		\1 $\m{A}$ é triangular $\implies$ $\det(\m{A})=a_{11} a_{22} \cdots a_{nn}$
		\1 $\m{A}$ é ortogonal $\implies$ $\det(\m{A}) = \pm1$
		\1 Seja $\m{B}$ resultante de uma das seguintes operações em linhas (colunas) de $\m{A}$.
			\2 Multiplicar linha (coluna) por $k$: $\det(\m{B}) = k\det(\m{A})$
			\2 Permutar linha (coluna): $\det(\m{B}) = -\det(\m{A})$
			\2 Múltiplo de uma linha (coluna) é somado a uma outra linha (coluna): $\det(\m{B}) = \det(\m{A})$
	\end{outline}
\end{theorem}

\section{Espaço vetorial Euclidiano}

\begin{definition} Seja $\v{v} \in \Re^n$.
	\begin{outline}[enumerate]
		\1 Norma: \( \norm{\v{v}} = \sqrt{v_1^2 + v_2^2 + \cdots + v_n^2} \)
	\end{outline}
\end{definition}

\begin{definition}[Produto escalar em $\Re^2$ e $\Re^3$] Sejam $\v{u}, \v{v} \in \Re^2 \text{ ou } \Re^3$ e $\theta$ o ângulo entre $\v{u}$ e $\v{v}$. O produto escalar entre $\v{u}$ e $\v{v}$ é definido como
	\begin{equation}
		\v{u} \cdot \v{v} = \norm{\v{u}} \norm{\v{v}} \cos{\theta}
	\end{equation}
\end{definition}

\begin{definition}[Produto escalar em $\Re^n$] Sejam $\v{u}, \v{v} \in \Re^n$. O produto escalar entre $\v{u}$ e $\v{v}$ (produto interno euclidiano) é definido como
	\begin{equation}
		\v{u} \cdot \v{v} = u_1 v_1 + u_2 v_2 + \dots + u_n v_n
	\end{equation}
	\begin{note}
		\leavevmode
		\begin{outline}[enumerate]
			\1 Em notação matricial: $\v{u} \cdot \v{v} = \v{u}\T \v{v} = \v{v}\T \v{u}$
			\1 No caso especial em que $\v{u} = \v{v}$, segue que 
			\begin{equation}
				\norm{\v{v}} = \sqrt{\v{u} \cdot \v{v}} = \sqrt{v_1^2 + v_2^2 + \cdots + v_n^2} 
			\end{equation}
		\end{outline}
	\end{note}
\end{definition}

\begin{definition} A distância entre $\v{u} \in \Re^n$ e $\v{v} \in \Re^n$ é definida por
	\begin{equation}
		d(\v{u}, \v{v}) = \norm{\v{u}-\v{v}} = \sqrt{(u_1-v_1)^2 + \cdots + (u_n-v_n)^2}
	\end{equation}
\end{definition}

\begin{definition}[Ortogonalidade de vetores] Sejam $\v{u}, \v{v} \in \Re^n$. $\v{u}, \v{v}$ são ortogonais se $\v{u} \cdot \v{v} = 0$.
\end{definition}

\begin{theorem}[Proriedades da norma e produto interno]
	\leavevmode
	\begin{outline}[enumerate]
		\1 \( \norm{\v{v}} \geq 0 \)
		\1 \( \norm{\v{v}} = 0 \iff \v{v} = \v{0} \)
		\1 \( \norm{k\v{v}} =\abs{k}\norm{\v{v}} \)
		\1 \( \abs{ \v{u} \cdot \v{v}} \leq \norm{\v{u}}\norm{\v{v}} \) (Desigualdade de Cauchy-Schwarz)
		\1 \( \norm{ \v{u} + \v{v}} \leq \norm{\v{u}} + \norm{\v{v}} \) (Desigualdade triangular)
		\1 \( \norm{ \v{u} + \v{v}}^2 = \norm{\v{u}}^2 + \norm{\v{v}}^2 \) (Teorema de Pitágoras em $\Re^n$)
		\1 \( d(\v{u}, \v{v}) \leq d(\v{u}, \v{2}) + d(\v{w}, \v{v})  \) (Desigualdade triangular para distâncias)
		\1 \( \norm{ \v{u} + \v{v}}^2 + \norm{ \v{u} - \v{v}}^2 = 2(\norm{\v{u}}^2 + \norm{\v{v}}^2) \) (Equações do paralelogramo para vetores)
		\1 \( \v{u} \cdot \v{v} = \frac{1}{4}\norm{ \v{u} + \v{v}}^2 + \frac{1}{4}\norm{ \v{u} - \v{v}}^2 \)
		\1 \(\m{A} \v{u} \cdot \v{v} = \v{u} \cdot \m{A}\T \v{v} \)
		\1 \(\v{u} \cdot \m{A} \v{v} = \m{A}\T  \v{u} \cdot \v{v} \)
	\end{outline}
\end{theorem}

\begin{theorem}[Teorema da projeção]
	Se $\v a$ e $\v u$ são vetores em $\Re^n$, e se $\v a \neq \v 0$, então $\v u$ pode ser expresso exatamente em uma única forma $\v u = \v w_1 + \v w_2$, onde $\v w_1$ é um múltiplo escalar de $\v a$ e $\v w_2$ é ortogonal a $\v a$. Além disso,
	\begin{equation}
		\v w_1 = \frac{\v u \cdot \v a}{\norm a^2}\v a \quad\quad \v w_2 = \v u - \v w_1
	\end{equation}
	\begin{note}
		\leavevmode
		\begin{outline}[enumerate]
			\1 $\v w_1$ é chamado de projeção ortogonal de $\v u$ em $\v a$ e denotado por $\proj_{\v a}{\v u}$;
			\1 $\v w_2$ é chamado de componente de $\v u$ ortogonal a $\v a$.
		\end{outline}
	\end{note}
\end{theorem}

\begin{theorem}
	\begin{equation}
		\norm{\proj_{\v a}{\v u}} = \norm{ \v u}  \abs {\cos(\theta)}
	\end{equation}
	%
	onde $\theta$ é o ângulo entre $\v u$ e $\v a$.
\end{theorem}

\section{Espaço vetorial Geral}

\begin{definition}[Combinação linear]
	Seja $\v w$ um vetor no espaço vetorial $V$. Dizemos que $\v w$ é uma combinação linear de vetores $\v v_1, \v v_2, \ldots, \v v_r$ em $V$ se $\v w$ pode ser escrito na forma
	\begin{equation}
		\v w = k_1 \v v_1 + k_2 \v v_2 + \ldots + k_r \v v_r
	\end{equation}
	%
	onde $k_1, k_2, \ldots, k_r$ são escalares chamados coeficientes da combinação linear.
\end{definition}

\begin{theorem} Seja $S = \{ \v w_1, \ldots, \v w_k \}$ um conjunto não vazio de vetores em $V$. Então:
	\begin{outline}[enumerate]
		\1 O conjunto $W$ de todas as possíveis combinações lineares de vetores em $S$ é um subespaço de $V$.
		\1 O cojunto $W$ (do item acima) é o menor subespaço de $V$ que contém todos os vetores de $S$, ou seja, qualquer outro subespaço contém todos os vetores de $W$.
	\end{outline}
\end{theorem}

\begin{definition}[Subespaço gerado]
	Seja $S = \{ \v w_1, \ldots, \v w_k \}$ um conjunto não vazio de vetores em $V$ e $W$ o subespaço de todas as possíveis combinações lineares dos vetores em $S$. Este subespaço é chamado de \textit{subespaço de $V$  gerado por $S$}, e dizemos que $\v w_1, \ldots, \v w_k$ geram $W$. Denotamos este subespaço por
	\begin{equation}
		W = \ger \{ \v w_1, \ldots, \v w_k \} \quad \text{ou} \quad  W = \ger(V)
	\end{equation}
\end{definition}

\begin{definition}[Conjunto LI]
	Seja $S = \{ \v v_1, \ldots, \v v_r \}$ um conjunto de dois ou mais vetores de um espaço vetorial $V$. Dizemos que $S$ é um conjunto linearmente independente (LI) se nenhum vetor em $S$ pode ser expresso como combinação linear de outros.
\end{definition}

\begin{theorem}
	Seja $S = \{ \v w_1, \ldots, \v w_k \} \subset V$ um conjunto não vazio. $V$ é LI se e somente se os coeficientes que satisfazem a equação 
	\begin{equation}
		k_1 \v r_1 + \cdots + k_r \v v_r = 0
	\end{equation}
	% 
	são $k_1=k_2=\cdots=k_r = 0$.
\end{theorem}

\begin{theorem}[4.3.2, Anton]
	\leavevmode
	\begin{outline}[enumerate]
		\1 Um conjunto finito que contém $\v{0}$ é linearmente dependente.
		\1 Um conjunto com exatamente um vetor é linearmente independente se e somente se esse vetor não é $\v{0}$.
		\1 Um conjunto com exatamente dois vetores é linearmente independente se e somente se nenhum dos dois vetores é múltiplo escalar do outro.
	\end{outline}
\end{theorem}

\begin{theorem}[4.3.3, Anton]
	Seja $S = \{\v{v}_1,\v{v}_2,\ldots,\v{v}_r\}$ um conjunto de vetores em $\Re^n$. Se $r>n$, então $\v{S}$ é linearmente dependente.
\end{theorem}

\begin{definition}[Base de espaço vetorial]
	Se $S = \{\v{v}_1,\v{v}_2,\ldots,\v{v}_n\}$ é um conjunto de vetores em um espaço vetorial $V$ de dimensão finita, então $\v{S}$ é chamado uma \textbf{base} para $V$ se:
	\begin{outline}[enumerate]
		\1 $S$ gera $V$.
		\1 $S$ é linearmente independente.
	\end{outline}
\end{definition}

\begin{theorem}[Unicidade da representação em base]
	Se $S = \{\v{v}_1,\v{v}_2,\ldots,\v{v}_n\}$ é uma base para um espaço vetorial $V$, então todo vetor $\v{v}$ em $V$ pode ser expresso na forma $\v{v} = c_1\v{v}_1 + c_2\v{v}_2 + \cdots + c_n\v{v}_n$ de exatamente uma forma.
\end{theorem}

\begin{theorem}[Teorema do Mais/Menos]
	Seja $S$ um conjunto não vazio de vetores em um espaço vetorial $V$.
	\begin{outline}[enumerate]
		\1 Se $S$ é um conjunto linearmente independente, e se $\v{v}$ é um vetor em $V$ que está fora de $\ger(S)$, então o conjunto $S\cup\{\v{v}\}$ que resulta ao inserir $\v{v}$ em $S$ ainda é linearmente independente.
		\1 Se $\v{v}$ é um vetor em $S$ que é expressável como uma combinação linear dos demais vetores em $S$, e se $S-\{\v{v}\}$ denota o conjunto obtido ao remover $\v{v}$ de $S$, então $S$ e $S-\{\v{v}\}$ geram o mesmo espaço; ou seja,
		\begin{equation}
			\ger(S) = \ger(S-\{\v{v}\})
		\end{equation}
	\end{outline}
\end{theorem}

\begin{theorem}[Dimensão de subespaço]
	Se $W$ é um subespaço de um espaço vetorial $V$ de dimensão finita, então:
	\begin{outline}[enumerate]
		\1 $W$ é de dimensão finita.
		\1 $\dim(W) \leq \dim(V)$.
		\1 $W = V$ se e somente se $\dim(W) = \dim(V)$.
	\end{outline}
\end{theorem}

\begin{theorem}[Base a partir de $n$ vetores]
	Seja $V$ um espaço vetorial $n$-dimensional, e seja $S$ um conjunto em $V$ com exatamente $n$ vetores. Então $S$ é uma base para $V$ se e somente se $S$ gera $V$ ou $S$ é linearmente independente.
\end{theorem}

\begin{theorem}[Matriz de transição para a base canônica]
	Seja $B' = \{\v{u}_1,\v{u}_2,\ldots,\v{u}_n\}$ uma base qualquer para o espaço vetorial $\Re^n$ e seja $S = \{\v{e}_1,\v{e}_2,\ldots,\v{e}_n\}$ a base canônica para $\Re^n$. Se os vetores dessas bases são escritos na forma de coluna, então
	\begin{equation}
		\m{P}_{B'\to S} = \left[\begin{array}{c|c|c|c}
			\v{u}_1 & \v{u}_2 & \cdots & \v{u}_n
		\end{array}\right]
	\end{equation}
\end{theorem}


\section{Espaços com produto interno}

\begin{definition}[Conjunto ortogonal e ortonormal]
	Um conjunto de dois ou mais vetores em um espaço com produto interno real é dito \textbf{ortogonal} se todos os pares de vetores distintos no conjunto são ortogonais. Um conjunto ortogonal no qual cada vetor tem norma 1 é dito \textbf{ortonormal}.
\end{definition}

\begin{theorem}[Independência linear de conjuntos ortogonais]
	Se $S = \{\v{v}_1,\v{v}_2,\ldots,\v{v}_n\}$ é um conjunto ortogonal de vetores não nulos em um espaço com produto interno, então $S$ é linearmente independente.
\end{theorem}


\begin{theorem}[Teorema da Projeção]
	Se $W$ é um subespaço de dimensão finita de um espaço com produto interno $V$, então todo vetor $\v{u}$ em $V$ pode ser expresso de exatamente uma forma como
	\begin{equation}
		\v{u} = \v{w}_1 + \v{w}_2
	\end{equation}
	em que $\v{w}_1$ está em $W$ e $\v{w}_2$ está em $W^{\perp}$.
	\begin{note}
		\leavevmode
		\begin{outline}[enumerate]
			\1 Os vetores $\v{w}_1$ e $\v{w}_2$ na equação acima são comumente denotados por
			\begin{equation}
				\v{w}_1 = \proj_W \v{u} \quad \text{e} \quad \v{w}_2 = \proj_{W^{\perp}} \v{u}
			\end{equation}
			Estes são chamados de \textit{projeção ortogonal de $\v{u}$ em $W$} e \textit{projeção ortogonal de $\v{u}$ em $W^{\perp}$}, respectivamente. O vetor $\v{w}_2$ também é chamado de \textit{componente de $\v{u}$ ortogonal a $W$}. Usando essa notação, a decomposição pode ser expressa como
			\begin{equation}
				\v{u} = \proj_W \v{u} + \proj_{W^{\perp}} \v{u}
			\end{equation}
		\end{outline}
	\end{note}
\end{theorem}

\begin{theorem}[Representação em base ortogonal/ortonormal e projeção sobre subespaço]
	\leavevmode
	\begin{outline}[enumerate]
		\1 Se $S = \{\v{v}_1,\v{v}_2,\ldots,\v{v}_n\}$ é uma base ortogonal para um espaço com produto interno $V$, e se $\v{u}$ é um vetor qualquer em $V$, então
		\begin{equation}
			\v{u} = \frac{\ip{\v{u}}{\v{v}_1}}{\norm{\v{v}_1}^2}\v{v}_1 + \frac{\ip{\v{u}}{\v{v}_2}}{\norm{\v{v}_2}^2}\v{v}_2 + \cdots + \frac{\ip{\v{u}}{\v{v}_n}}{\norm{\v{v}_n}^2}\v{v}_n
		\end{equation}
		\1 Se $S = \{\v{v}_1,\v{v}_2,\ldots,\v{v}_n\}$ é uma base ortonormal para um espaço com produto interno $V$, e se $\v{u}$ é um vetor qualquer em $V$, então
		\begin{equation}
			\v{u} = \ip{\v{u}}{\v{v}_1}\v{v}_1 + \ip{\v{u}}{\v{v}_2}\v{v}_2 + \cdots + \ip{\v{u}}{\v{v}_n}\v{v}_n
		\end{equation}
		\1 Seja $W$ um subespaço de dimensão finita de um espaço com produto interno $V$. Se $\{\v{v}_1,\v{v}_2,\ldots,\v{v}_r\}$ é uma base ortogonal para $W$, e $\v{u}$ é um vetor qualquer em $V$, então
		\begin{equation}
			\proj_W \v{u} = \frac{\ip{\v{u}}{\v{v}_1}}{\norm{\v{v}_1}^2}\v{v}_1 + \frac{\ip{\v{u}}{\v{v}_2}}{\norm{\v{v}_2}^2}\v{v}_2 + \cdots + \frac{\ip{\v{u}}{\v{v}_r}}{\norm{\v{v}_r}^2}\v{v}_r
		\end{equation}
		\1 Se $\{\v{v}_1,\v{v}_2,\ldots,\v{v}_r\}$ é uma base ortonormal para $W$, e $\v{u}$ é um vetor qualquer em $V$, então
		\begin{equation}
			\proj_W \v{u} = \ip{\v{u}}{\v{v}_1}\v{v}_1 + \ip{\v{u}}{\v{v}_2}\v{v}_2 + \cdots + \ip{\v{u}}{\v{v}_r}\v{v}_r
		\end{equation}
	\end{outline}
	\begin{note}
		\leavevmode
		\begin{outline}[enumerate]
			\1 Usando a terminologia e notação da Definição 2 da Seção 4.4, segue dos itens 1 e 2 que o vetor de coordenadas de um vetor $\v{u}$ em $V$ em relação a uma base ortogonal $S = \{\v{v}_1,\v{v}_2,\ldots,\v{v}_n\}$ é
			\begin{equation}
				(\v{u})_S = \left(\frac{\ip{\v{u}}{\v{v}_1}}{\norm{\v{v}_1}^2}, \frac{\ip{\v{u}}{\v{v}_2}}{\norm{\v{v}_2}^2}, \ldots, \frac{\ip{\v{u}}{\v{v}_n}}{\norm{\v{v}_n}^2}\right)
			\end{equation}
			\1 e em relação a uma base ortonormal $S = \{\v{v}_1,\v{v}_2,\ldots,\v{v}_n\}$ é
			\begin{equation}
				(\v{u})_S = \left(\ip{\v{u}}{\v{v}_1}, \ip{\v{u}}{\v{v}_2}, \ldots, \ip{\v{u}}{\v{v}_n}\right)
			\end{equation}
		\end{outline}
	\end{note}
\end{theorem}


\begin{theorem}[Decomposição espectral]
	Se $\m{A}$ é uma matriz simétrica diagonalizada ortogonalmente por
	\begin{equation}
		\m{P} = \left[\v{u}_1 \ \ \v{u}_2 \ \ \cdots \ \ \v{u}_n\right],
	\end{equation}
	e se $\lambda_1,\lambda_2,\ldots,\lambda_n$ são os autovalores de $\m{A}$ correspondentes aos autovetores unitários $\v{u}_1,\v{u}_2,\ldots,\v{u}_n$, então $\m{A}$ pode ser expressa como
	\begin{equation}
		\m{A} = \sum_{i=1}^n \lambda_i \v{u}_i\v{u}_i\T.
	\end{equation}
\end{theorem}

\section{Posto e nulidade}

\begin{definition}[Posto e nulidade] Seja $\m{A}$ qualquer;
	\leavevmode
	\begin{outline}[enumerate]
		\1 Posto de $\m{A}$ ($\pos(\m{A})$): dimensão comum do espaço-linha e do espaço-coluna de $\m{A}$;
		\1 Nulidade de $\m{A}$ ($\nul(\m{A})$): dimensão do espaço-nulo de $\m{A}$.
	\end{outline}
\end{definition}

\begin{theorem}[Espaços nulo, linha e coluna]
	\leavevmode
	\begin{outline}[enumerate]
		\1 Operações elementares em linhas não alteram o espaço-nulo da matriz;
		\1 Operações elementares em linhas não alteram o espaço-linha da matriz. \uline{Porém, podem alterar o espaço-coluna};
		\1 Se $\m{A}$ está na forma escalonada por linhas (\textit{row echelon form}), então os vetores linha com líderes $1$ formam uma base para o espaço-linha de $\m{A}$, e os vetores coluna com os líderes $1$ formam uma base para o espaço-coluna de $\m{A}$; 
		\1 Sejam $\m{A}$ e $\m{B}$ matrizes equivalentes por linhas. Então:
			\2 Um dado conjunto de vetores colunas de $\m{A}$ é LI $\iff$ os vetores colunas correspondentes em $\m{B}$ são LI.
			\2 Um dado conjunto de vetores colunas de $\m{A}$ forma uma base para oe espaço-coluna de $\m{A}$ $\iff$ os vetores colunas correspondentes em $\m{B}$ formam uma base para o espaço-coluna de $\m{B}$.
	\end{outline}
\end{theorem}

\begin{theorem}[Propriedades do posto e nulidade] Seja $\m{A}_{m\times n}$
	\leavevmode
	\begin{outline}[enumerate]
		\1 \( \pos(\m{A}) = \pos(\m{A}\T) \)
		\1 \( \pos(\m{A}) = \pos(\m{A}\T\m{A}) =  \pos(\m{A}\m{A}\T)\)
		\1 \( \pos(\m{A}) + \nul(\m{A}) = n \)
		\1 $\pos(\m{A})$: núm. de variáveis líderes na solução geral de $\m{A}\v{x}=\v{b}$
		\1 $\nul(\m{A})$: núm. de parâmetros na solução geral de $\m{A}\v{x}=\v{b}$
		\1 Se $\pos(\m{A}) = p$, então $n-p$ raízes do polinômio característico são nulos 
	\end{outline}
\end{theorem}

\begin{theorem}[Teorema da Consistência] Sja $\m{A}\v{x}=\v{b}$ é um sistema linear de $m$ equações por $n$ incógnitas \underline{com $\v{b}$ fixado.} As seguintes afirmações são equivalentes:
	\leavevmode
	\begin{outline}[enumerate]
		\1 $\m{A}\v{x}=\v{b}$ é consistente.
		\1 $\v{b}$ está no espaço-coluna de $\m{A}$.
		\1 $\m{A}$ e $[\m{A} \mid \v{b}]$ tem o mesmo posto. 
	\end{outline}
\end{theorem}

\begin{theorem}Seja $\m{A}\v{x}=\v{b}$ um sistema linear de $m$ equações por $n$ incógnitas. As seguintes afirmações são equivalentes:
	\leavevmode
	\begin{outline}[enumerate]
		\1 $\m{A}\v{x}=\v{b}$ é consistente para todo $\v{b}$.
		\1 Os vetores-coluna $\m{A}$ geram o $\Re^m$.
		\1 $\pos(\m{A})=m$. 
	\end{outline}
\end{theorem}

\begin{theorem}Seja $\m{A}\v{x}=\v{b}$ um sistema linear \underline{consistente} de $m$ equações por $n$ incógnitas. Se $\pos(\m{A})=r$, então a solução geral do sistema contém $n-r$ parâmeros.
\end{theorem}

\begin{theorem}[Sistemas $m \times n$] Seja $\m{A}_{m\times n}$. As seguintes afirmações são equivalentes:
	\leavevmode
	\begin{outline}[enumerate]
		\1 $\m{A}\v{x}=\v{0}$ possui somente a solução trivial.
		\1 Os vetores-coluna de $\m{A}$ são LI.
		\1 $\m{A}\v{x}=\v{b}$ tem no máximo uma (uma ou nenhuma)solução para cada $\v{b}$.
	\end{outline}
\end{theorem}





\section{Autovalores e autovetores}

\begin{definition} Definições sobre autovalores e autovetores.
	\leavevmode
	\begin{outline}[enumerate]
		\1 Seja $\m{A}_{n\times n}$. Um vetor não nulo $\v{x}\in \Re^n$ é chamado de     \textbf{autovetor} de $\m{A}$ se 
			\begin{equation}
				\m{A} \v{x} = \lambda \v{x}
			\end{equation}
		% 
		para algum escalar $\lambda$, o qual é chamado de \textbf{autovalor} de $\m{A}$.
		\1 Equação característico: $\det(\lambda \m{I} - \m{A})=0$
		\1 Quando a equação característica é expandida, ela toma a forma
			\begin{equation}
				p(\lambda) = \lambda^n + c_1\lambda^{n-1} + \cdots + c_n
			\end{equation}
		%
		e $p(\lambda)$ é denominado polinômio característico.
		\1 Espaço-solução do sistema $(\lambda \m{I} - \m{A})\v{x}=0$: \textbf{auto-espaço} de $\m{A}$ associado a $\lambda$.
		\1 Multiplicidade algébrica de $\lambda_0$: número de vezes que $(\lambda - \lambda_0)$ aparece como um fator do polinômio característico.
		\1 Multiplicidade geométrica de $\lambda_0$: dimensão do auto-espaço associado a $\lambda-\lambda_0$
	\end{outline}
\end{definition}

\begin{theorem}[Propriedades de autovalores e autovetores] Seja $\m{A}_{n\times n}$.
	\leavevmode
	\begin{outline}[enumerate]
		\1 $\m{A}$ é triangular $\implies$ os autovalores de $\m{A}$ são as entradas da diagonal principal
		\1 \( \m{A}^k \v{x}= \lambda^k\v{x} \)
		\1 $\m{A}$ é invertível $\iff$ $\lambda=0$ não é um autovalor de $\m{A}$
	\end{outline}
\end{theorem}


\begin{theorem}[Propriedades das raízes característica] Seja $\m{A}_{n\times n}$.
	\leavevmode
	\begin{outline}[enumerate]
		\1 A soma das raízes carecterísticas de $\m{A}$ é $\tr(\m{A})$
		\1 O produto das raízes carecterísticas de $\m{A}$ é $\det(\m{A})$
		\1 Se $\pos(\m{A}) = p$, então $n-p$ raízes do polinômio característico são nulos 
	\end{outline}
\end{theorem}


\section{Semelhança e Diagonalização}

\begin{definition}[Matrizes semelhantes]
	Sejam $\m{A}$ e $\m{B}$ quadradas. $\m{B}$ é semelhante a $\m{A}$ se existe uma matriz $\m{P}$ tal que $\m{B} = \m{P}\inv \m{A}\m{P}$.
	\begin{note}
		\begin{outline}[enumerate]
			Se $\m{B}$ é similar a $\m{A}$, então $\m{A}$ é similar $\m{B}$.
		\end{outline}
	\end{note}
\end{definition}

\begin{theorem}[Propriedades de matrizes semelhantes]
	Sejam $\m{A}$ e $\m{B}=\m{P}\inv\m{A} P$ semelhantes. As seguintes propriedades de $\m{A}$ são preservadas em $\m{B}$:
	\begin{outline}[enumerate]
		\1 Determinante
		\1 Invertibilidade
		\1 Posto
		\1 Nulidade
		\1 Traço
		\1 Polinômio característico
		\1 Autovalores
		\1 Dimensões dos autoespaços.
	\end{outline}
\end{theorem}

\begin{definition}
	$\m{A}$ quadrada é dita diagonalizável se é semelhante a alguma matriz diagonal, isto é, se existe uma matriz invertível $\m{P}$ tal que $\m{P}\inv A \m{P}$ é diagonal. Neste caso, dizemos que a matriz $\m{P}$ diagonaliza $\m{A}$.
\end{definition}

\begin{theorem}[Diagonalização] Seja $\m{A}_{n\times n}$.
	\leavevmode
	\begin{outline}[enumerate]
		\1 $\m{A}$ é diagonalizável $\iff$ $\m{A}$ tem $n$ autovetores LI.
		\1 $\m{A}$ é diagonalizável $\iff$ Para cada autovalor, a multiplicidade geométrica é igual à multiplicidade algébrica.
		\1 Se $\lambda_1, \ldots, \lambda_k$ são autovalores distintos de $\m{A}$ e $\v{v}_1, \ldots, \v{v}_k$ são os autovetores correspondentes, então $\{\v{v}_1, \ldots, \v{v}_k\}$ é um conjunto LI.
		\1 $\m{A}$ tem $n$ autovalores distintos $\implies$ $\m{A}$ é diagonalizável.
		\1 $\m{A}$ é diagonalizável $\implies \m{A}^k = \m{P}\m{D}^k\m{P}\inv$.
	\end{outline} 
\end{theorem}




\section{Diagonalização ortogonal}

\begin{theorem}[Matrizes ortogonais e bases ortonormais] Seja $\m{A}_{n \times n}$. As seguintes afirmações são equivalentes.
	\begin{outline}[enumerate]
		\1 $\m{A}$ é ortogonal.
		\1 os vetores linha de $\m{A}$ formam uma base ortonormal em $\Re^n$ com produto interno euclidiano.
		\1 os vetores coluna de $\m{A}$ formam uma base ortonormal em $\Re^n$ com produto interno euclidiano.
		\1 \( \norm{\m{A}\v{x}} = \norm{\v{x}} \quad \forall \v{x}\in \Re^n \)
		\1 \( \m{A}\v{x}\cdot\m{A}\v{y} = \v{x}\cdot\v{y} \quad \forall \v{x},\v{y} \in \Re^n \)
	\end{outline}
\end{theorem}

\begin{theorem}[Propriedades de matrizes ortogonais] Seja $\m{A}_{n \times n}$.
	\begin{outline}[enumerate]
		\1 $\m{A}$ é ortogonal $\implies \m{A}\T$ é ortogonal.
		\1 $\m{A}$ é ortogonal $\implies \m{A}\inv$ é ortogonal.
		\1 $\m{A}, \m{B}$ são ortogonais $\implies \m{A}\m{B}$ é ortogonal.
		\1 $\m{A}$ é ortogonal $\implies \det(\m{A})=\pm 1$
	\end{outline}
\end{theorem}

\begin{theorem}[Mudança de bases ortonormais] Seja $V$ um espaço finito-dimensional com produto interno. Se $\m{P}$ é a matriz de transição de uma base ortonormal de $V$ para uma outra base ortonormal para $V$, então $\m{P}$ é ortogonal.
\end{theorem}

\begin{definition}[Semelhança e diagonalização ortogonal]
	\leavevmode
	\begin{outline}[enumerate]
		\1 Se $\m{A}, \m{B}$ são matrizes quadradas, dizemos que $\m{B}$ é \uline{ortogonalmente semelhante} a $\m{A}$ se existe uma matriz ortogonal $\m{P}$ tal que $\m{B}=\m{P}\T\m{A}\m{P}$.
		\1 Se $\m{B}$ é ortogonalmente semelhante a alguma matriz diagonal $\m{D}$, ou seja,
		\begin{equation}
			\m{P}\T\m{A}\m{P} = \m{D},
		\end{equation}
		% 
		dizemos que $\m{A}$ é ortogonalmente diagonalizável e que $\m{P}$ diagnoliza (ortogonalmente) $\m{A}$. 
	\end{outline}
	\begin{note}
		\begin{outline}[enumerate]
			\1 Se $\m{B}$ é ortogonalmente semelhante a $\m{A}$, então $\m{A}$ é ortogonalmente semelhante a $\m{B}$.
		\end{outline}
	\end{note}
\end{definition}

\begin{theorem}[Diagonalização ortogonal] As seguintes afirmações são equivalentes.
	\begin{outline}[enumerate]
		\1 $\m{A}$ é ortogonalmente diagonalizável.
		\1 $\m{A}$ tem um conjunto ortonormal de $n$ autovetores.
		\1 $\m{A}$ é simétrica.
	\end{outline}
\end{theorem}


\begin{theorem}[Matrizes simétricas] Seja $\m{A}_{n\times n}$ simétrica com entradas reais. Então:
	\leavevmode
	\begin{outline}[enumerate]
		\1 Os autovalores de $\m{A}$ são todos reais.
		\1 Autovetores de autoespaços diferentes são ortogonais.
	\end{outline} 
\end{theorem}

\section{Formas lineares e quadráticas}


\begin{definition}[Formas lineares e quadráticas]
	\leavevmode
	\begin{outline}[enumerate]
		\1 Forma linear: $b_1 x_1 + b_2 x_2 + \cdots + b_nx_n$
			\2 $f(\v{x}) = \v{b}\T\v{x} = \v{x}\T\v{b}$
		\1 Forma afim: $f(\v{x}) = a + \v{b}\T\v{x}$
		\1 Forma quadrática: $a_1 x_1^2 + a_2 x_2^2 + \cdots + a_nx_n^2$ (mais todos possíveis termos $a_kx_ix_j$, com $i\neq j$)
			\2 Se $\m{A}_{n\times n}$ é simétrica e $\v{x}$ é vetor-coluna $n\times 1$, a forma quadrática associada a $\m{A}$ é $Q_\m{A} = \v{x}\T\m{A}\v{x}$
	\end{outline}
	\begin{note}
		\leavevmode
		\begin{outline}[enumerate]
			\1 Se $\m{A}$ é a diagonal, a forma quadrática não possui termos cruzados.
		\end{outline}
		conteúdo...
	\end{note}
\end{definition}


\begin{definition}[Classificação de formas quadráticas e matrizes] Uma forma quadrática $\v{x}\T\m{A}\v{x}$ e a respectiva matriz $\m{A}$ simétrica é dita:
	\begin{outline}[enumerate]
		\1 Positiva definida: se $\v{x}\T\m{A}\v{x} > 0$ para todo $\v{x} \neq 0$ 
		\1 Positiva semidefinida: se $\v{x}\T\m{A}\v{x} \geq 0$ para todo $\v{x} \neq 0$ e  $\v{x}\T\m{A}\v{x} = 0$ para pelo menos um $\v{x} \neq 0$
		\1 Negativa definida: se $\v{x}\T\m{A}\v{x} < 0$ para todo $\v{x} \neq 0$ 
		\1 Negativa semidefinida: se $\v{x}\T\m{A}\v{x} \leq 0$ para todo $\v{x} \neq 0$ e  $\v{x}\T\m{A}\v{x} = 0$ para pelo menos um $\v{x} \neq 0$
		\1 Indefinida: $\v{x}\T\m{A}\v{x}$ assume valores positivos e negativos.
	\end{outline}
\end{definition}

\begin{theorem} 
	Seja $\m{A}_{n\times n}$ simétrica.
	\begin{outline}[enumerate]
		\1 $\v{x}\T\m{A}\v{x}$ ($\m{A}$) é positiva definida $\iff$ todos os autovalores de $\m{A}$ são positivos.
		\1 $\v{x}\T\m{A}\v{x}$ ($\m{A}$) é negativa definida $\iff$ todos os autovalores de $\m{A}$ são negativos.
		\1 $\v{x}\T\m{A}\v{x}$ ($\m{A}$) é indefinido $\iff$ $\m{A}$ tem pelo menos um autovalor positivo e pelo menos um autovalor negativo.
	\end{outline}
\end{theorem}

\begin{theorem} 
	Seja $\m{A}_{n\times n}$ simétrica. As seguintes afirmações são equivalentes.
	\begin{outline}[enumerate]
		\1 $\m{A}$ é positiva definida 
		\1 Todos os autovalores de $\m{A}$ são positivos.
		\1 Os determinantes de cada submatriz da diagonal principal são positivos (critério de Sylvester).
	\end{outline}
\end{theorem}

\begin{theorem} Seja $\m{A}_{n\times n}$ simétrica. Se $\m{A}$ é positiva definida, então:
	\begin{outline}[enumerate]
		\1 $\pos(\m{A})=n$
		\1 $a_{ii}>0$ para todo $i=1,\ldots,n$
		\1 $\m{P}\T\m{A}\m{P}$ é positiva definida para toda matriz $\m{P}_{n \times p}$
	\end{outline}
\end{theorem}

\begin{theorem} Seja $\m{A}_{n\times n}$ simétrica. Se $\m{A}$ é positiva semidefinida, então:
	\begin{outline}[enumerate]
		\1 $\pos(\m{A})<n$
		\1 $a_{ii}\geq0$ para todo $i=1,\ldots,n$
		\1 $\m{P}\T\m{A}\m{P}$ é positiva semidefinida para toda matriz $\m{P}_{n \times p}$
	\end{outline}
\end{theorem}

\begin{theorem} 
	Seja $\m{A}_{m \times n}$. Então, 
	\begin{outline}[enumerate]
		\1 $\m{A}\T\m{A}$ é positiva semidefinida. 
		\1 $\m{A}\T\m{A}$ é positiva se tem posto-coluna completo, ou seja, $\pos(\m{A})$. 
	\end{outline}
\end{theorem}


\section{Matrizes em blocos}

\begin{theorem}
	Seja $\m{B}$ uma matriz particionada em blocos,
	\begin{equation}
		\m{B} =
		\begin{bmatrix}
			\m{B}_{11} & \m{B}_{12} \\
			\m{B}_{21} & \m{B}_{22}
		\end{bmatrix},
	\end{equation}
	em que $\m{B}_{11}$ e $\m{B}_{22}$ são matrizes quadradas \underline{não-singulares}. Então
	\begin{align}
		\m{B}\inv &=
		\begin{bmatrix}
			(\m{B}_{11} - \m{B}_{12}\m{B}_{22}\inv \m{B}_{21})\inv & -\m{B}_{11}\inv \m{B}_{12}(\m{B}_{22} - \m{B}_{21}\m{B}_{11}\inv \m{B}_{12})\inv \\[2ex]
			-\m{B}_{22}\inv \m{B}_{21}(\m{B}_{11} - \m{B}_{12}\m{B}_{22}\inv \m{B}_{21})\inv & (\m{B}_{22} - \m{B}_{21}\m{B}_{11}\inv \m{B}_{12})\inv
		\end{bmatrix}
	\end{align}
	e
	\begin{equation}
		\deter{\m{B}} = \deter{\m{B}_{22}}\deter{\m{B}_{11} - \m{B}_{12}\m{B}_{22}\inv \m{B}_{21}} = \deter{\m{B}_{11}}\deter{\m{B}_{22} - \m{B}_{21}\m{B}_{11}\inv \m{B}_{12}}.
	\end{equation}
\end{theorem}

\begin{theorem}
	Seja $\m{B}$ uma matriz bloco diagonal,
	\begin{equation}
		\m{B} =
		\begin{bmatrix}
			\m{B}_{11} & \m{0} & \m{0} \\
			\m{0} & \m{B}_{22} & \m{0} \\
			\m{0} & \m{0} & \m{B}_{33}
		\end{bmatrix},
	\end{equation}
	em que $\m{B}_{11}$, $\m{B}_{22}$ e $\m{B}_{33}$ são matrizes quadradas \underline{não-singulares}. Então
	\begin{equation}
		\m{B}\inv =
		\begin{bmatrix}
			\m{B}_{11}\inv & \m{0} & \m{0} \\
			\m{0} & \m{B}_{22}\inv & \m{0} \\
			\m{0} & \m{0} & \m{B}_{33}\inv
		\end{bmatrix}
	\end{equation}
	e
	\begin{equation}
		\deter{\m{B}} = \deter{\m{B}_{11}}\deter{\m{B}_{22}}\deter{\m{B}_{33}}.
	\end{equation}
	O resultado vale para matrizes bloco diagonal de maior dimensão.
\end{theorem}

\section{Derivadas}

\begin{definition}[Derivada de função escalar]
	Considere uma função escalar $f : \Re^n \longrightarrow \Re$, em que $\v{b}\in\Re^n$. Adotaremos como convenção a derivada da função $f$ em relação ao vetor $\v{b}$ como o gradiente
	\begin{equation}
		\frac{\partial f}{\partial \v{b}} = \nabla_{\v{b}} f =
		\begin{pmatrix}
			\dfrac{\partial f}{\partial b_1} \\[1.5ex]
			\dfrac{\partial f}{\partial b_2} \\[0.5ex]
			\vdots \\[0.5ex]
			\dfrac{\partial f}{\partial b_n}
		\end{pmatrix}.
	\end{equation}
	Assim, $\dfrac{\partial f}{\partial \v{b}}$ é um vetor de dimensão $n\times 1$.
\end{definition}

\begin{theorem}
	Considere a forma linear $f(\v{x}) = \v{b}\T\v{x} = \v{x}\T\v{b}$, em que $\v{b}\in\Re^n$ é um vetor constante. Como $f(\v{x}) = b_1x_1 + b_2x_2 + \cdots + b_nx_n$, temos
	\begin{equation}
		\frac{\partial f}{\partial x_i} = b_i, \quad i=1,\ldots,n.
	\end{equation}
	Portanto,
	\begin{equation}
		\boxed{\dfrac{\partial}{\partial \v{x}}\left(\v{b}\T\v{x}\right) = \dfrac{\partial}{\partial \v{x}}\left(\v{x}\T\v{b}\right) = \v{b}}.
	\end{equation}
\end{theorem}

\begin{theorem}
	Considere $f(\v{x}) = \v{x}\T\m{A}\v{x}$. A derivada em relação a $x_k$ é
	\begin{equation}
		\frac{\partial f}{\partial x_k} = \underbrace{\sum_{j=1}^n a_{kj}x_j}_{\text{linha } k \text{ de } \m{A}} + \underbrace{\sum_{i=1}^n a_{ik}x_i}_{\text{coluna } k \text{ de } \m{A}}.
	\end{equation}
	A derivada em $\v{x}$ nesse caso é dada por
	\begin{equation}
		\boxed{\frac{\partial}{\partial \v{x}}\left(\v{x}\T\m{A}\v{x}\right) = (\m{A}+\m{A}\T)\v{x}.}
	\end{equation}
	Em particular, se $\m{A}$ é simétrica (forma quadrática), então $\m{A}\T=\m{A}$, logo $(\m{A}+\m{A}\T)\v{x} = 2\m{A}\v{x}$ e, portanto,
	\begin{equation}
		\boxed{\frac{\partial}{\partial \v{x}}\left(\v{x}\T\m{A}\v{x}\right) = 2\m{A}\v{x}.}
	\end{equation}
\end{theorem}

\begin{theorem}
	Considere $f(\v{x}) = \v{x}\T\m{A}\v{x} + \v{b}\T\v{x}$. Então
	\begin{equation}
		\frac{\partial f}{\partial \v{x}} = (\m{A}+\m{A}\T)\v{x} + \v{b}.
	\end{equation}
	Se $\m{A}$ é simétrica,
	\begin{equation}
		\boxed{\frac{\partial f}{\partial \v{x}} = 2\m{A}\v{x} + \v{b}.}
	\end{equation}
\end{theorem}